\usepackage{bm}
\usepackage{tikz}

\theoremstyle{theorem}
\newtheorem*{thmA}{Theorem A}
\newtheorem*{thmB}{Theorem B}
\newtheorem*{thmC}{Theorem C}
\newtheorem*{thmD}{Theorem D}
\newtheorem*{thmE}{Theorem E}
\newtheorem*{conjecture*}{Conjecture}
\newtheorem*{lemmaA}{Lemma}

\newcommand{\shrink}[2]{#1_{\left>#2\right<}}
\newcommand{\dcara}{\dd_C}
\newcommand{\pql}{$\psi$-ql }
\newcommand{\toto}{\rightrightarrows}

\newcommand{\BA}{\mathbf{A}}
\newcommand{\FB}{\mathfrak{B}}
\newcommand{\fl}{\mathfrak{l}}
\newcommand{\FL}{\mathfrak{L}}
\newcommand{\FD}{\mathfrak{D}}

\newcommand{\Mol}{\MM ol}

\newcommand{\tot}{{\mathrm{tot}}}
\newcommand{\inner}{{\mathrm{inn}}}
\newcommand{\comp}{{\mathrm{comp}}}
\newcommand{\biv}{{\mathrm{biv}}}
\newcommand{\br}{{\mathrm{br}}}
\newcommand{\nbr}{{\mathrm{ubr}}}
\newcommand{\ver}{{\mathrm{ver}}}
\newcommand{\hor}{{\mathrm{hor}}}
\newcommand{\bd}{{\mathrm{bd}}}
\newcommand{\nbd}{{\mathrm{nbd}}}
\newcommand{\fun}{{\mathrm{fun}}}
\newcommand{\str}{{\mathrm{str}}}
\newcommand{\sta}{{\mathrm{sta}}}
\newcommand{\cra}{{\mathrm{cra}}}
\newcommand{\std}{{\mathrm{std}}}
\newcommand{\cen}{{\mathrm{cen}}}
\newcommand{\rip}{{\mathrm{rip}}}
\newcommand{\van}{{\mathrm{van}}}
\newcommand{\out}{\operatorname{out}}
\newcommand{\bgg}{\mathbf G}
\newcommand{\inbr}{{\mathrm{in}}}
\newcommand{\splitwave}{{\mathrm{split}}}
\newcommand{\overwave}{{\mathrm{over}}}
\newcommand{\outbr}{{\mathrm{out}}}
\newcommand{\amp}{{\mathrm{amp}}}

\newcommand{\etf}{{E(T_f)}}
\newcommand{\rt}[1]{\operatorname{rt}(#1)}




\newcommand{\wcan}{W_{\operatorname{can}}}
\newcommand{\fcan}{\FF_{\operatorname{can}}}
\newcommand{\ren}[1]{\RR #1}

\newcommand{\dominates}{\multimap}
\newcommand{\weakarrows}{\rightarrow}
\newcommand{\weakdominates}{\dominates}
\newcommand{\bh}{\mathbf H}
\newcommand{\hperp}{\HH^\biv}
\newcommand{\hv}{\HH^\ver}
\newcommand{\topdom}{\rightsquigarrow}


\newcommand{\p}{\mathbf p}
\newcommand{\bound}{\mathbf b}
\newcommand{\n}{\mathbf n}
\newcommand{\q}{\mathbf q}
\newcommand{\wstd}{\WW(\RR F)}
\newcommand{\Flux}{\operatorname{Flux}_{\operatorname{std}}}



\theoremstyle{definition}
\numberwithin{equation}{section}
\numberwithin{figure}{section}




\newcommand{\com} {{ \mathrm{com}}}
\newcommand{\thin} {{\mathrm{thin}}}
\newcommand{\thick} {{ \mathrm{thick}}}
\newcommand{\loc}{{\mathrm {loc}}}
\newcommand{\ess}{{\mathrm{ess}}}
\newcommand{\ness}{{\mathrm{ne}}}
\newcommand{\pro}{{\mathrm{pr}}}

\newcommand{\locflux} {\WW^\perp_\loc}
\DeclareMathOperator{\flux}{flux}
\newcommand{\stdflux}[1]{\flux(#1)}
\newcommand{\ddi}{\dd^i}

%\newcommand{\re}{{\mathrm{(r)}}}
\newcommand{\essre}{{\mathrm{ess(r)}}}
\newcommand{\nere}{{\mathrm{ne(r)}}}
\newcommand{\verre}{{\mathrm{ver(r)}}}

\newcommand{\Cyl}{{\mathrm{Cyl}}}
\newcommand{\hsubset}{{ \,\underset{h}{\subset}\, }    }
\newcommand{\hsupset}{{ \, \underset{h}{\supset} \, }     }

% \newcommand{\lp}{\multimap}
\newcommand{\lp}{{\,\mathrm {dominates}\,}}

\DeclareMathOperator{\disk}{disk}
\DeclareMathOperator{\preper}{preper}
\DeclareMathOperator{\val}{Val}

\newcommand{\clone}[3]{#1^{#2}_{#3}}
\newcommand{\PF}{Perron-Fr\"obenius}

\newcommand{\FJ}{\mathfrak{J}}
\newcommand{\BU}{\mathbf{U}}
\newcommand{\FK}{\mathfrak{K}}
\newcommand{\FT}{\mathfrak{T}}
\newcommand{\BT}{\mathbf{T}}
\newcommand{\BK}{\mathbf{K}}
\newcommand{\per}{\mathbf{n}}
\newcommand{\translationregion}{T}

%\input{ak_macros}

\makeatletter
\newtheorem{repeatlem@}{Lemma}
\newenvironment{repeatlem}[1]{%
	\def\therepeatlem@{\ref{#1}}
	\repeatlem@
}
{\endrepeatlem@}
\makeatother

\makeatletter
\newtheorem{repeatprop@}{Proposition}
\newenvironment{repeatprop}[1]{%
	\def\therepeatprop@{\ref{#1}}
	\repeatprop@
}
{\endrepeatprop@}
\makeatother


\newcommand{\good}{^g}
\newcommand{\bad}{^b}
\newcommand{\vg}{^{g\dominates}}
\newcommand{\gdom}{\vg}
\newcommand{\gb}{^{gb\dominates}}
\newcommand{\bv}{^{b}}
\newcommand{\gbg}{^{gb\dominates g}}
\newcommand{\gdomg}{^{g \dominates g}}

\newcommand{\xwbv}{X}
%\newcommand{\xhv}{X^{h+v}}
\newcommand{\xhv}{X}

\newcommand{\jj}{\jmath}

\newcommand{\Card}{Card}

% Define a command for a cardioid symbol
% The 'scale' factor controls the overall size. 
% For a symbol roughly the size of an M, a scale of about 0.05 to 0.1 might be appropriate 
% for a standard 10pt or 12pt article font size.
\newcommand{\cardioidsymbol}{%
    \begin{tikzpicture}[scale=0.08, baseline=(current bounding box.center)]
        \draw[line width=0.5pt, smooth] 
            plot[domain=0:360, samples=360, variable=\t] 
            ({2*(1+cos(\t))*cos(\t)}, {2*(1+cos(\t))*sin(\t)});
    \end{tikzpicture}%
}

\renewcommand{\Card}{\cardioidsymbol}

% Define a new command for the adjusted cardioid symbol
% New modifications:
% 1. Belly button to the right: Achieved by setting the polar equation domain from -180 to 180 (or 0 to 360) 
%    and changing the 'x' and 'y' coordinates in the plot.
% 2. Lift it up three points: Achieved with the 'yshift=3pt' option.
% 3. Shrink it 15\%: Achieved by reducing the scale factor by 15\%. 
%    Since the previous scale was 0.08, the new scale is 0.08 * (1 - 0.15) = 0.068.
\newcommand{\cardioidsymbolRightShifted}{%
    \begin{tikzpicture}[
        scale=0.068, % Reduced 15\% from the original 0.08 scale
        yshift=3pt, % Lifted up 3 points
        baseline=(current bounding box.center)
    ]
        % Polar equation: r = 2(1 + cos(t))
        % Standard Cartesian plot: (r*cos(t), r*sin(t))
        % For a right-facing cardioid, the standard plot is used over the full domain.
        \draw[line width=0.5pt, smooth] 
            plot[domain=0:360, samples=360, variable=\t] 
            ({2*(1+cos(\t))*cos(\t)}, {2*(1+cos(\t))*sin(\t)});
    \end{tikzpicture}%
}

\renewcommand{\Card}{\cardioidsymbolRightShifted}

\renewcommand{\Card}{\operatorname{Card}}

\newcommand{\gc}[1]{\breve{#1}}

\newcommand{\tunedby}{\bot}

\newcommand{\gr}[1]{{\color{Green} #1}}

\newcommand{\alex}[1]{{\color{Cyan} #1}}


\newcommand{\bK}{\mathbf K}
\newcommand{\mathhyphen}{\text{-}}





